\setmainhead{\emph{\textsc{Ilíada}}, Canto IX}
\bnum
\section*{\centering\textbf{CANTO IX}\footT{\ehdos, pp. 73-74.}}
\addcontentsline{toc}{subsection}{\textsc{Canto IX}}

\separator
\postart
\setotherline{29}
Dijo, y todos inmóviles y en silencio han quedado.\\
Largo tiempo afligida calló la gente aquea\\
hasta que al fin Diomedes el buen jefe así ha hablado:

---Atrida, rechazo ante todo tu absurda idea,\\
pues tal es mi derecho de rey en la asamblea.\\
Mas, no te enojes, aunque mi honra heriste primero,\\
dándome entre los dánaos por flojo y mal guerrero;\\
lo que, de todo argivo, mozo y viejo, es sabido.\\
Si el hijo del artero Kronos, entre dos cosas,\\
darte el honor supremo del cetro ha decidido,\\
te negó el valor que es la fuerza más poderosa.\\
¡Desdichado! ¿Creíste que es tan floja y de escaso\\
valor la gente aquea como dices? Si acaso\\
al regreso te incita tu ánimo, parte ya,\\
que abierta senda tienes y junto al mar está\\
la gran flota que desde Micenas te siguier.\\
Mas los otros aqueos de larga cabellera,\\
hasta que a Troya hundamos se quedarán. Y si\\
también al caro suelo patrio huyen navegando,\\
yo y Esténelo, solos, quedaremos peleando\\
Hasta el fin de Ilión, puesto que dios nos trajo aquí.\par
\poend
\separator
%
\postart
\setotherline{308}
Divino Laertiades, Ulises ingenioso\\
preciso es que os exponga con veraz desembozo\\
lo que yo pienso y cómo lo he de cumplir, y dónde,\\
para que ceséis vuestro circunloquio enfadoso.\\
Pues, igual que las puertas del Hades, me es odioso\\
quien dice lo contrario de lo que en su alma esconde.\par
\poend
\separator
%
\postart
\setotherline{340}
¿Acaso, entre los hombres, sólo aman los Atridas\\
a sus esposas? Todo noble y cuerdo varón,\\
quiere y cuida la suya, cual yo, de corazón\\
amé a la que tuve, aunque por las armas fué habida.\par
\poend
\separator
%
\postart
\setotherline{697}
Agamenón rey de hombres, gloriosísimo Atrida,\\
nunca al irreprochable Peliónida debiste\\
suplicar, prometiéndole donación tan crecida;\\
pues si ya era orgulloso, mucho más lo engreiste.\\
Por lo que hace a él, dejémoslo que se vaya o se quede.\\
Cuando le dé la gana, de nuevo peleará,\\
o cuando un dios lo incite; pero lo que procede,\\
yo os lo diré y lo haremos: id a acostaros ya,\\
después que hayan saciado vuestro apetito, ahora,\\
comida y vino, que esto fuerza y valor nos dá;\\
y al despuntar la bella rosodáctila Aurora,\\
forma tú ante las naves gente y carros ligeros,\\
y a todos animando, lucha entre los primeros.
\poend
\enum